\documentclass[11pt]{article}
\usepackage[a4paper,margin=1in]{geometry}
\usepackage{amsmath,amssymb}
\usepackage[hidelinks]{hyperref}
\title{A Counterexample to Conjecture 00000008558}
\author{gaochengzhecpu (AI-assisted submission)}
\date{}
\setlength{\emergencystretch}{3em}
\begin{document}
\maketitle

The conjectured equality between Boolean sublattice embedding dimension
and the minimum number of lattice generators is false, even for the
three-element distributive chain
\[
C_3=\{0<a<1\}.
\]
Its embedding dimension is $2$.  Its minimum number of generators is
$3$ when the lattice operations are $\wedge,\vee$, and is $1$ when
the bounds $0,1$ are also available as constants.  Thus both standard
signature conventions contradict the proposed equality.

\paragraph{Embedding dimension is exactly two.}
For a nonnegative integer $n$, write $B_n=\mathcal P(\{0,\ldots,n-1\})$,
with intersection and union as meet and join.  A sublattice embedding
is an injective map preserving both operations.
Define
\[
f(0)=\varnothing,\qquad f(a)=\{0\},\qquad f(1)=\{0,1\}.
\]
These are three distinct, nested subsets.  On a chain, meet and join
are minimum and maximum; on these nested subsets, intersection and
union are the corresponding smaller and larger subsets.  Therefore
$f:C_3\hookrightarrow B_2$ is a sublattice embedding.
It also preserves bottom and top, so the construction works if that
additional convention is imposed.

The lattices $B_0$ and $B_1$ have respectively one and two elements,
so neither admits an injection from $C_3$.  Consequently the minimum
embedding dimension exists and equals $2$.
This counterexample lies entirely within the class of lattices which
do embed into a finite Boolean lattice.

\paragraph{Generators with the two binary lattice operations.}
The usual lattice signature has precisely $\wedge$ and $\vee$;
the bounded-lattice signature additionally includes two constants
\cite{burris}.  First use the ordinary lattice signature.
For any $x,y\in C_3$, both $x\wedge y$ and $x\vee y$ belong to
$\{x,y\}$.  Induction on lattice terms therefore shows that every
term whose inputs lie in a set $S\subseteq C_3$ evaluates to an
element already in $S$.  Hence the generated set is exactly $S$.
To generate $C_3$, all three elements must be chosen, and the full
set does generate it.  The minimum number is thus $3\ne2$.
Generation here means closure under both lattice operations, as in
the source's statement about generators of a lattice.

\paragraph{The convention with bounds as constants.}
If $0$ and $1$ may be used without being counted as generators,
the same term induction shows that $S$ generates exactly
$S\cup\{0,1\}$.  A generating set must contain $a$, and $\{a\}$
does generate the whole chain.  The minimum number is therefore
$1\ne2$.  This covers the other standard convention without
changing the embedding dimension.

\paragraph{Formal verification and semantic map.}
The accompanying Lean 4.19.0 file imports only \texttt{Std}.
The three elements are represented by \texttt{Fin 3}, with numerical
labels $0,1,2$ for the bottom, middle, and top respectively.
The order, minimum, and maximum satisfy all lattice laws and both
distributive identities.

\texttt{BooleanLattice n} is the full type of Boolean membership
functions on \texttt{Fin n}.  Intersection and union are pointwise
Boolean conjunction and disjunction.  Arbitrary proposition-valued
subsets have a proved membership-function encoding.
\texttt{BooleanEmbedding} requires injectivity and preservation of
both operations.  The file constructs the two-dimensional embedding
and proves that it also preserves bounds.  It proves that no embedding
in dimensions $0$ or $1$ exists, quantifying over all maps, rather
than merely checking a chosen list of candidate embeddings.
\texttt{boolean\_dimension\_two} establishes existence and minimality;
\texttt{boolean\_dimension\_unique} identifies every minimum as $2$.

Every possible generating set, including every proposition-valued
subset, is encoded by its three membership bits.  The predicates
\texttt{Generated} and \texttt{BoundedGenerated} are inductively
defined by the actual lattice operations, with constants only in
the latter.  Their characterization is proved by induction, so it
covers terms of arbitrary depth.  The least-closure universal property
is also proved for ordinary generation.
The two minimum-generator predicates assert both attainment and a
lower bound over every generating set.  Their minima are formally
proved to be $3$ and $1$ respectively.

The final theorems
\begin{quote}
\texttt{conjecture\_8558\_false}\\
\texttt{conjecture\_8558\_bounded\_signature\_false}
\end{quote}
negate the existence of a common minimum for dimension and generator
number under the two conventions.  All three minima have separately
been proved to exist, so the disagreement is between actual finite
numbers.  This disproves the second conjunct of the source conjecture;
the separate proposed bound involving join-irreducible elements is
not needed.  Neither theorem uses an admitted proof, a native proof
oracle, or an additional axiom.  Both final theorems depend only on
\texttt{propext} and \texttt{Quot.sound}.

\begin{thebibliography}{1}
\bibitem{burris}
S. Burris and H. P. Sankappanavar,
\emph{A Course in Universal Algebra}, Millennium Edition, 2012 update,
Chapter II, Section 1, examples (8)--(9), and Section 3.
\url{https://www.math.uwaterloo.ca/~snburris/htdocs/UALG/univ-algebra2012.pdf}
\end{thebibliography}
\end{document}
